% GENERATED FILE. Edit canonical lessons, then run npm run generate:books.
% canonical-source-hash: fnv1a64:801233e3a14f8e91
% canonical-lessons: LA-C124-codex, LA-C124-atramentum, LA-C124-calamus, LA-C124-cera, LA-C124-pictura

\chapter{Book, Ink, Reed pen, Wax, Picture}
\label{ch:la-book-ink-reed-pen-wax-picture}
\hlchaptermodality{124}

\begin{chapteropening}
\textbf{By the end of this chapter:} \emph{I can say a book, ink, a reed pen, wax and a picture.}

Complete the last lesson of chapter 124: i can say a book, ink, a reed pen, wax and a picture.
\end{chapteropening}

\section[cōdex, cōdicis]{cōdex, cōdicis — a book (of bound pages)}

\label{lesson:LA-C124-codex}



\begin{quote}
Before the new one: say the Latin for a purse, then the Latin for a coin.

Read the account aloud as one run: \emph{Māne surgō. In scholā legō et scrībō. Vespere domum redeō.}
\end{quote}

\subsection*{You'll want to know: cōdex, cōdicis}
\textbf{cōdex, cōdicis} — \textquotedblleft{}a book (of bound pages)\textquotedblright{}.

\begin{grammarlens}[title={What you've built}]
One more word for this chapter.
\end{grammarlens}

\subsection*{Your turn}
\begin{itemize}
\raggedright
  \item \emph{Say it:} \emph{cōdex}
  \item \emph{Say it:} \emph{cōdex}, once more
  \item \emph{Say it:} say \emph{cōdex}
  \item \emph{Recall:} say the Latin for a purse, then the Latin for a coin, then say \emph{cōdex} again
\end{itemize}

\begin{tcolorbox}[breakable,colback=teal!4,colframe=teal!35!black,title={Before you move on}]
What does cōdex, cōdicis mean? (\textquotedblleft{}A book (of bound pages)\textquotedblright{}.) Say it once more.
\end{tcolorbox}
\section[ātrāmentum, ātrāmentī]{ātrāmentum, ātrāmentī — ink}

\label{lesson:LA-C124-atramentum}



\begin{quote}
Before the new one: say the Latin for a coin, then the Latin for a book.
\end{quote}

\subsection*{You'll want to know: ātrāmentum, ātrāmentī}
\textbf{ātrāmentum, ātrāmentī} — \textquotedblleft{}ink\textquotedblright{}.

\begin{grammarlens}[title={What you've built}]
One more word for this chapter.
\end{grammarlens}

\subsection*{Your turn}
\begin{itemize}
\raggedright
  \item \emph{Say it:} \emph{ātrāmentum}
  \item \emph{Say it:} \emph{ātrāmentum}, once more
  \item \emph{Say it:} say \emph{ātrāmentum}
  \item \emph{Recall:} say the Latin for a coin, then the Latin for a book, then say \emph{ātrāmentum} again
\end{itemize}

\begin{tcolorbox}[breakable,colback=teal!4,colframe=teal!35!black,title={Before you move on}]
What does ātrāmentum, ātrāmentī mean? (\textquotedblleft{}Ink\textquotedblright{}.) Say it once more.
\end{tcolorbox}
\section[calamus, calamī]{calamus, calamī — a reed pen}

\label{lesson:LA-C124-calamus}



\begin{quote}
Before the new one: say the Latin for a book, then the Latin for ink.
\end{quote}

\subsection*{You'll want to know: calamus, calamī}
\textbf{calamus, calamī} — \textquotedblleft{}a reed pen\textquotedblright{}.

\begin{grammarlens}[title={What you've built}]
One more word for this chapter.
\end{grammarlens}

\subsection*{Your turn}
\begin{itemize}
\raggedright
  \item \emph{Say it:} \emph{calamus}
  \item \emph{Say it:} \emph{calamus}, once more
  \item \emph{Say it:} say \emph{calamus}
  \item \emph{Recall:} say the Latin for a book, then the Latin for ink, then say \emph{calamus} again
\end{itemize}

\begin{tcolorbox}[breakable,colback=teal!4,colframe=teal!35!black,title={Before you move on}]
What does calamus, calamī mean? (\textquotedblleft{}A reed pen\textquotedblright{}.) Say it once more.
\end{tcolorbox}
\section[cēra, cērae]{cēra, cērae — wax; a wax tablet}

\label{lesson:LA-C124-cera}



\begin{quote}
Before the new one: say the Latin for ink, then the Latin for a reed pen.
\end{quote}

\subsection*{You'll want to know: cēra, cērae}
\textbf{cēra, cērae} — \textquotedblleft{}wax; a wax tablet\textquotedblright{}.

\begin{grammarlens}[title={What you've built}]
One more word for this chapter.
\end{grammarlens}

\subsection*{Your turn}
\begin{itemize}
\raggedright
  \item \emph{Say it:} \emph{cēra}
  \item \emph{Say it:} \emph{cēra}, once more
  \item \emph{Say it:} say \emph{cēra}
  \item \emph{Recall:} say the Latin for ink, then the Latin for a reed pen, then say \emph{cēra} again
\end{itemize}

\begin{tcolorbox}[breakable,colback=teal!4,colframe=teal!35!black,title={Before you move on}]
What does cēra, cērae mean? (\textquotedblleft{}Wax; a wax tablet\textquotedblright{}.) Say it once more.
\end{tcolorbox}
\section[pictūra, pictūrae]{pictūra, pictūrae — a picture, a painting}

\label{lesson:LA-C124-pictura}



\begin{quote}
Before the new one: say the Latin for a reed pen, then the Latin for wax; a wax tablet.
\end{quote}

\subsection*{You'll want to know: pictūra, pictūrae}
\textbf{pictūra, pictūrae} — \textquotedblleft{}a picture, a painting\textquotedblright{}.

\begin{grammarlens}[title={What you've built}]
One more word for this chapter.
\end{grammarlens}

\subsection*{Your turn}
\begin{itemize}
\raggedright
  \item \emph{Say it:} \emph{pictūra}
  \item \emph{Say it:} \emph{pictūra}, once more
  \item \emph{Say it:} say \emph{pictūra}
  \item \emph{Recall:} say the Latin for a reed pen, then the Latin for wax; a wax tablet, then say \emph{pictūra} again
\end{itemize}

\begin{tcolorbox}[breakable,colback=teal!4,colframe=teal!35!black,title={Before you move on}]
What does pictūra, pictūrae mean? (\textquotedblleft{}A picture, a painting\textquotedblright{}.) Say it once more.
\end{tcolorbox}
